% !TeX program = xelatex
\documentclass[11pt,a4paper]{ctexart}
\usepackage[color]{tlnote}
\title{高数 I（A 类）· 薄弱点定向 网络题源精练}
\author{陶理}
\noteheader{48 题 · 六模块 · 对标交大 2008–2025 十七套}
\docdate{2026-11-10}
\begin{document}
\renewcommand{\arraystretch}{1.15}
\noindent{\Large\bfseries 薄弱点定向 · 网络题源精练（48 题）}\hfill{\small\color{tlblue}2026-11-10}
\par\noindent\rule{\linewidth}{0.4pt}\par\smallskip
\noindent{\small\color{gray}针对六模块薄弱点，从公开课件/题库检索题型，按上海交大 2008–2025 十七套期中真题难度重设数值（与库内 84 题去重）。\textbf{题目为对标重设，不是交大原题}；全部答案经 sympy 独立复核。}
\begin{keypoint}[使用说明]
\begin{itemize}\itemsep2pt
\item 每模块 8 题：基础 2 / 进阶 4 / 压轴 2。先做后对答案，卡住满 5 min 再看提示。
\item 答案速查在册末；错题按「按错因回炉索引」回对应章节。
\item 难度对齐：作图 $\approx$ 真题倒数第二题 11 分；量词 $\approx$ 第 11 题 8 分；导数块 16–20 分；泰勒大题 8 分。
\end{itemize}
\end{keypoint}
\section{M1 函数性态作图七步}
\noindent{\small\color{tlblue}A 级第 1 名 · 约 11 分/16 套 · 倒数第二题}\\[0.2em]
\noindent{\footnotesize\color{gray}题源参考：北师大《微积分》Ch3 课件（math0.bnu.edu.cn/~liujun，作图七步例 13）、复旦《函数图象描绘》讲义（math.fudan.edu.cn）、MIT OCW 18.01SC Ses28b。均按交大真题难度重设函数。}\\[0.5em]
\qi{M1.1}{基础·斜+垂直}{求 $y=\dfrac{x^{2}+2}{x-1}$ 的斜渐近线与垂直渐近线。}
\qi{M1.2}{进阶·斜+垂直}{求 $y=\dfrac{x^{3}}{x^{2}-4}$ 的斜渐近线与垂直渐近线。}
\qi{M1.3}{进阶·极值+渐近}{求 $y=x+\dfrac1x$ 的垂直渐近线、斜渐近线与极值点。}
\qi{M1.4}{进阶·水平+拐点}{求 $y=x\ee^{-x}$ 的水平渐近线、极值点与拐点。}
\qi{M1.5}{压轴·极值+拐点}{求 $y=x^{2}\ee^{-x}$ 的水平渐近线、极值点与拐点横坐标。}
\qi{M1.6}{基础·水平+极值}{求 $y=\dfrac{x^{2}}{x^{2}-1}$ 的垂直渐近线、水平渐近线与极值点。}
\qi{M1.7}{进阶·两端斜渐近}{求 $y=\sqrt{x^{2}+1}$ 在 $x\to+\infty$ 与 $x\to-\infty$ 的两条渐近线。}
\qi{M1.8}{压轴·两端斜渐近}{求 $y=x\arctan x$ 的两条斜渐近线。}
\section{M2 量词行、$\eps$-$\delta$ 与 M 型}
\noindent{\small\color{tlblue}A 级第 2 名 · 17/17 套必考 8 分 · 第 11 题}\\[0.2em]
\noindent{\footnotesize\color{gray}题源参考：Emory 大学 $\eps$-$\delta$ 练习（math.oxford.emory.edu/site/math111/probSetEpsilonDelta）、STEM-Master 极限定义题、Pearson 教材习题。按交大第 11 题三种形式重设。}\\[0.5em]
\qi{M2.1}{基础·$\eps$-$\delta$}{用定义证 $\lim\limits_{x\to3}(4x-1)=11$，给出 $\delta(\eps)$。}
\qi{M2.2}{进阶·先关区间}{用定义证 $\lim\limits_{x\to2}(x^{2}+1)=5$，给出 $\delta(\eps)$。}
\qi{M2.3}{进阶·先关区间}{用定义证 $\lim\limits_{x\to1}(2x^{2}-x)=1$，给出 $\delta(\eps)$。}
\qi{M2.4}{基础·有界×无穷小}{用定义证 $\lim\limits_{x\to0}x\cos\dfrac1x=0$，给出 $\delta(\eps)$。}
\qi{M2.5}{进阶·$\eps$-$X$}{用定义证 $\lim\limits_{x\to+\infty}\dfrac{x+3}{x-1}=1$，给出 $X(\eps)$。}
\qi{M2.6}{基础·$\eps$-$X$}{用定义证 $\lim\limits_{x\to+\infty}\dfrac1{x^{2}}=0$，给出 $X(\eps)$。}
\qi{M2.7}{进阶·M 型}{用定义证 $\lim\limits_{x\to0}\dfrac1{x^{4}}=+\infty$，给出 $\delta(M)$。}
\qi{M2.8}{进阶·M 型·非零趋近点}{用定义证 $\lim\limits_{x\to2}\dfrac1{(x-2)^{2}}=+\infty$，给出 $\delta(M)$。}
\section{M3 条件书写与六型求导}
\noindent{\small\color{tlblue}A 级第 3 名 · 导数块 16→20 分}\\[0.2em]
\noindent{\footnotesize\color{gray}题源参考：台大 OCW 高阶导/隐函数（ocw.aca.ntu.edu.tw）、同济《高等数学》习题 2-4（参数/隐函数）、大学搜题酱参数+隐函数二阶题。按交大导数块分值重设。}\\[0.5em]
\qi{M3.1}{进阶·参数二阶}{$x=t^{2}+1,\ y=t^{3}-t$，求 $\dfrac{\D y}{\D x},\dfrac{\D^{2}y}{\D x^{2}}$ 在 $t=1$ 处的值；写前提。}
\qi{M3.2}{进阶·参数二阶}{$x=\ee^{t}\cos t,\ y=\ee^{t}\sin t$，求 $\dfrac{\D y}{\D x},\dfrac{\D^{2}y}{\D x^{2}}$ 在 $t=0$ 处的值；写前提。}
\qi{M3.3}{进阶·参数·化简}{$x=t^{2},\ y=t+\ln t\ (t>0)$，求 $\dfrac{\D y}{\D x}\big|_{t=2}$；写前提。}
\qi{M3.4}{基础·隐函数}{$x^{2}y+y^{3}=4$ 在 $(1,1)$ 处求 $y'$；写前提。}
\qi{M3.5}{进阶·隐函数·指数}{$\ee^{y}+xy=2$ 在 $(0,\ln2)$ 处求 $y'$；写前提。}
\qi{M3.6}{基础·分段可导}{$f(x)=\begin{cases}x^{2}+a,&x\le1\\ bx,&x>1\end{cases}$ 在 $x=1$ 可导，求 $a,b$。}
\qi{M3.7}{压轴·反函数二阶}{$f(x)=x^{3}+2x$，$y=f^{-1}(x)$，求 $y''(0)$；写前提。}
\qi{M3.8}{进阶·对数求导}{$y=x^{x}\ (x>0)$，求 $y'(1)$；写前提。}
\section{M4 泰勒：展开阶、系数与符号}
\noindent{\small\color{tlblue}A 级第 4 名 · 2023 起稳定大题 8 分}\\[0.2em]
\noindent{\footnotesize\color{gray}题源参考：浙大《极限计算 2：泰勒与洛必达》讲义（yuanhongyi.github.io/zjucalc25）、UC3M Calculus I Unit 8、常见泰勒极限题集。按交大泰勒大题重设。}\\[0.5em]
\qi{M4.1}{基础·通项系数}{求 $f(x)=(1+x)^{1/3}$ 的 $f'''(0)$。}
\qi{M4.2}{基础·展错阶}{求 $\lim\limits_{x\to0}\dfrac{\ee^{x}-1-x-\tfrac12x^{2}}{x^{3}}$。}
\qi{M4.3}{进阶·展错阶·符号}{求 $\lim\limits_{x\to0}\dfrac{\cos x-1+\tfrac12x^{2}-\tfrac1{24}x^{4}}{x^{6}}$。}
\qi{M4.4}{基础·展错阶}{求 $\lim\limits_{x\to0}\dfrac{\ln(1+x)-x+\tfrac12x^{2}}{x^{3}}$。}
\qi{M4.5}{压轴·高阶展错}{求 $\lim\limits_{x\to0}\dfrac{\tan x-x-\tfrac13x^{3}}{x^{5}}$。}
\qi{M4.6}{进阶·复合高阶导}{求 $f(x)=\ln(1+x^{2})$ 的 $f^{(6)}(0)$。}
\qi{M4.7}{压轴·幂指·泰勒}{求 $\lim\limits_{x\to0}\left(\dfrac1{\sin^{2}x}-\dfrac1{x^{2}}\right)$。}
\qi{M4.8}{进阶·级数乘移位}{求 $f(x)=x^{2}\ee^{-x}$ 的 $f^{(5)}(0)$。}
\section{M5 极限计算与间断点}
\noindent{\small\color{tlblue}B 级 · 间断 10/17 套}\\[0.2em]
\noindent{\footnotesize\color{gray}题源参考：清华社《高等数学》Ch1（tup.tsinghua.edu.cn）、Hanspub《隐秘单侧极限》论文、腾讯云考研极限与连续专题。按交大选填难度重设。}\\[0.5em]
\qi{M5.1}{基础·加减禁替换}{求 $\lim\limits_{x\to0}\dfrac{x-\tan x}{x^{3}}$。}
\qi{M5.2}{进阶·数列}{求 $\lim\limits_{n\to\infty}n\left(\sqrt{n^{2}+1}-n\right)$。}
\qi{M5.3}{进阶·根式}{求 $\lim\limits_{x\to+\infty}\left(\sqrt{x^{2}+x}-x\right)$。}
\qi{M5.4}{进阶·幂指}{求 $\lim\limits_{x\to0}(\cos x)^{1/x^{2}}$。}
\qi{M5.5}{基础·可去间断}{$f(x)=(1+x)^{1/x}$ 在 $x=0$ 无定义，判 $x=0$ 的间断类型。}
\qi{M5.6}{进阶·跳跃间断}{$f(x)=\arctan\dfrac1x$ 在 $x=0$ 的左右极限，并判型。}
\qi{M5.7}{进阶·对称展开}{求 $\lim\limits_{x\to0}\dfrac{\ee^{x}-\ee^{-x}-2x}{x^{3}}$。}
\qi{M5.8}{压轴·反三角对照}{求 $\lim\limits_{x\to0}\dfrac{\arcsin x-\sin x}{x^{3}}$。}
\section{M6 新压轴、函数方程与极坐标}
\noindent{\small\color{tlblue}B 级 · 2024-25 压轴转向 · 极坐标近 3 年连出}\\[0.2em]
\noindent{\footnotesize\color{gray}题源参考：台大 Ch10 参数方程与极坐标（math.ntu.edu.tw/~hchu）、科罗拉多 10.4 Polar Calculus、滑铁卢 MATH148 Ch4。按交大压轴与选填重设。}\\[0.5em]
\qi{M6.1}{进阶·极坐标切线}{极坐标 $r=1-\cos\theta$，求 $\theta=\dfrac\pi2$ 处切线斜率。}
\qi{M6.2}{进阶·极坐标切线}{极坐标 $r=\sin2\theta$，求 $\theta=\dfrac\pi4$ 处切线斜率。}
\qi{M6.3}{进阶·极坐标切线}{极坐标 $r=\ee^{-\theta}$，求 $\theta=0$ 处切线斜率。}
\qi{M6.4}{基础·函数方程}{$f(x+y)=f(x)f(y)$，$f$ 连续且 $f(1)=2$，求 $f(3)$。}
\qi{M6.5}{进阶·压缩迭代}{$a_1=2,\ a_{n+1}=\sqrt{2+a_n}$，求 $\lim\limits_{n\to\infty}a_n$。}
\qi{M6.6}{进阶·极坐标切线}{极坐标 $r=1+\sin\theta$，求 $\theta=\dfrac\pi3$ 处切线斜率。}
\qi{M6.7}{压轴·迭代·渐近阶}{$a_1=1,\ a_{n+1}=\dfrac{a_n}{1+a_n}$，求 $\lim\limits_{n\to\infty}na_n$。}
\qi{M6.8}{压轴·极坐标切线}{极坐标 $r=2\cos\theta$，求 $\theta=\dfrac\pi4$ 处切线斜率。}
\section{答案速查（48 行，只给结果）}
{\footnotesize
\begin{longtable}{@{}p{1.4cm}p{7.4cm}p{7.0cm}@{}}
\toprule
\textbf{题号} & \textbf{答案} & \textbf{卡点（一句话）} \\
\midrule
\endfirsthead
\toprule
\textbf{题号} & \textbf{答案} & \textbf{卡点（一句话）} \\
\midrule
\endhead
\bottomrule
\endlastfoot
M1.1 & 斜 $y=x+1$；垂直 $x=1$ & 多项式除法得 $b$，分母零点得垂直 \\
M1.2 & 斜 $y=x$；垂直 $x=\pm2$ & 分子比分母高一次，除后余项趋零 \\
M1.3 & 垂直 $x=0$；斜 $y=x$；极小 $(1,2)$，极大 $(-1,-2)$ & 两端都要查；$y'=1-1/x^{2}$ \\
M1.4 & 水平 $y=0$（仅 $x\to+\infty$）；极大 $\left(1,\tfrac1\ee\right)$；拐点 $\left(2,\tfrac{2}{\ee^{2}}\right)$ & 只有 $+\infty$ 一侧趋 0 \\
M1.5 & 水平 $y=0$；极小 $(0,0)$，极大 $\left(2,\tfrac{4}{\ee^{2}}\right)$；拐点 $x=2\pm\sqrt2$ & $y''=\ee^{-x}(x^{2}-4x+2)$ \\
M1.6 & 垂直 $x=\pm1$；水平 $y=1$；极大 $(0,0)$ & 同次看首项系数比 \\
M1.7 & $y=x$ 与 $y=-x$ & $\sqrt{x^{2}}=|x|$，两端符号不同 \\
M1.8 & $y=\tfrac{\pi}{2}x-1$ 与 $y=-\tfrac{\pi}{2}x-1$ & 两端斜率 $\pm\pi/2$，截距都是 $-1$ \\
M2.1 & $\delta=\eps/4$ & 先把 $|f-A|$ 化成 $4|x-3|$ \\
M2.2 & $\delta=\min\{1,\eps/5\}$ & 先限 $|x-2|<1$ 得 $|x+2|<5$ \\
M2.3 & $\delta=\min\{\tfrac12,\eps/4\}$ & $|2x+1|<4$ 当 $|x-1|<1/2$ \\
M2.4 & $\delta=\eps$ & $|\cos(1/x)|\le1$，别去求 $\cos$ 的极限 \\
M2.5 & $X=\max\{2,1+4/\eps\}$ & $|f-1|=\dfrac{4}{|x-1|}$ \\
M2.6 & $X=1/\sqrt{\eps}$ & 量词行别写成 $\delta$ \\
M2.7 & $\delta=M^{-1/4}$ & 要 $>M$：分母给下界 \\
M2.8 & $\delta=M^{-1/2}$ & 趋近点是 $2$ 不是 $0$ \\
M3.1 & $1$；$1$；前提 $x'_t=2t\ne0$ & 二阶要再除一次 $x'_t$ \\
M3.2 & $1$；$2$；前提 $x'_t=\ee^{t}(\cos t-\sin t)\ne0$ & 先化 $\tfrac{\sin+\cos}{\cos-\sin}$ \\
M3.3 & $\dfrac38$；前提 $t>0,\ x'_t=2t\ne0$ & 先化简再代值 \\
M3.4 & $-\dfrac12$；前提 $F_y=x^{2}+3y^{2}\ne0$ & 先验点在曲线上 \\
M3.5 & $-\dfrac{\ln2}{2}$；前提 $F_y=\ee^{y}+x\ne0$ & $F_y=\ee^{y}+x$，别漏 $x$ \\
M3.6 & $a=1,\ b=2$ & 先连续（$1+a=b$）再可导（$2=b$） \\
M3.7 & $0$；前提 $f'(0)=2\ne0$ & $y''=-\dfrac{f''(y)}{(f'(y))^{3}}$，$f''(0)=0$ \\
M3.8 & $1$；前提 $x>0$ & $\ln y=x\ln x$，乘回 $y$ \\
M4.1 & $\dfrac{10}{27}$ & 通项系数 $\alpha(\alpha-1)(\alpha-2)$，不乘 $3!$ \\
M4.2 & $\dfrac16$ & $\ee^{x}$ 展到三阶 \\
M4.3 & $-\dfrac1{720}$ & 分母六阶，$\cos x$ 展到六阶 \\
M4.4 & $\dfrac13$ & $\ln(1+x)$ 展到三阶 \\
M4.5 & $\dfrac{2}{15}$ & $\tan x$ 展到五阶 \\
M4.6 & $240$ & $\ln(1+u)$ 的三次项 $\tfrac{u^3}{3}$，乘 $6!$ \\
M4.7 & $\dfrac13$ & $\sin^{2}x$ 展到四阶 \\
M4.8 & $-20$ & $x^{5}$ 项系数 $-\tfrac1{3!}$，乘 $5!$ \\
M5.1 & $-\dfrac13$ & 加减式必须展开，不能逐项替换 \\
M5.2 & $\dfrac12$ & 有理化：分子分母同乘共轭 \\
M5.3 & $\dfrac12$ & 有理化后同除 $x$ \\
M5.4 & $\ee^{-1/2}$ & 取对数，$\ln\cos x\sim-\tfrac{x^{2}}2$ \\
M5.5 & 可去间断；补充定义 $f(0)=\ee$ & 极限存在但函数值缺 \\
M5.6 & 左 $-\tfrac\pi2$、右 $\tfrac\pi2$；跳跃间断 & $x\to0^{\pm}$ 时 $1/x\to\pm\infty$ \\
M5.7 & $\dfrac13$ & $\ee^{x}-\ee^{-x}=2\sinh x$ 展开 \\
M5.8 & $\dfrac13$ & 两式各展到三阶再相减 \\
M6.1 & $-1$ & 一律转参数式再求 $\D y/\D x$ \\
M6.2 & $-1$ & $r'(\pi/4)=0$，斜率由 $r$ 与角决定 \\
M6.3 & $-1$ & 参数式在 0 点：$x'=-1,\ y'=1$ \\
M6.4 & $8$ & 连续 + 乘法型 $\Rightarrow f(x)=2^{x}$ \\
M6.5 & $2$ & 解 $L=\sqrt{2+L}$，取非负根 \\
M6.6 & $-1$ & 心脏线中段斜率为 $-1$ \\
M6.7 & $1$ & 取倒数化等差：$1/a_n=n$ \\
M6.8 & $0$ & 该点为水平切线 \\
\end{longtable}
}
\section{按错因回炉索引}
\begin{itemize}\itemsep2pt
\item \textbf{渐近线漏 $b$ / 两端符号} $\to$ M1.1--M1.8；回精编本第 7 章七步卡。
\item \textbf{放缩方向反 / 量词写错} $\to$ M2.1--M2.8；回《极限定义证明》三模板。
\item \textbf{漏写前提 / 二阶导少除一次} $\to$ M3.1--M3.8；回冲刺计划 §5-C。
\item \textbf{展错阶 / 系数符号} $\to$ M4.1--M4.8；回精编本第 5 章。
\item \textbf{加减误用等价无穷小 / 左右极限} $\to$ M5.1--M5.8；回《函数极限与间断点》§4--§5。
\item \textbf{函数方程 / 迭代 / 极坐标} $\to$ M6.1--M6.8；回合订本 p206--p212。
\end{itemize}
\vfill
\noindent{\footnotesize 题源检索自公开课件与题库，按交大十七套真题难度重设数值并独立复核；非交大原题，不用于冒充真题。复核脚本：\texttt{build\_web\_drill.py}（全 PASS 才成文）。}
\end{document}